\documentclass{article}
\usepackage{graphicx}
\usepackage{amsmath}
\usepackage{geometry}

\geometry{margin=1in}

\title{Project and Demo Inventory}
\author{Your Name}
\date{\today}

\begin{document}

\maketitle

\section*{Projects and Demos}

\subsection*{Embedded Common C and C++}
\begin{itemize}
    \item \textbf{Static heap allocators}: Dynamic allocation semantics for advanced data structures such as open address hash tables, RB trees, and graphs.
    \item Compatible with \texttt{std::c++}.
    \item C++14 conforming \texttt{stevemac::vector} with static allocator.
\end{itemize}

\subsection*{STM32CubeIDE}
\begin{itemize}
    \item Interrupt bare metal "multi-threading" with \texttt{std::atomics}.
    \item System Graph.
    \item STM32CubeMX pinout configuration.
    \item Semihosting.
\end{itemize}

\subsection*{Keil uVision with STM32CubeMX}
\begin{itemize}
    \item STM32CubeMX pinout config integrated.
    \item Interrupt bare metal "multi-threading" with C11 atomics.
\end{itemize}

\subsection*{VSCode}
\begin{itemize}
    \item Rust embedded bare metal state captures.
    \item Semihosting.
    \item Cortex Quick Start for STM32F401RE.
\end{itemize}

\subsection*{Visual Studio}
\begin{itemize}
    \item Raspberry Pi remote cross compile and debugging.
\end{itemize}

\subsection*{MLibs}
\begin{itemize}
    \item Embedded automation suite.
    \item Semihosting.
\end{itemize}

\subsection*{Raspberry Pi Assembly}
\begin{itemize}
    \item Native assembly development (sorts and other routines).
\end{itemize}

\subsection*{Algorithms C++ General}
\begin{itemize}
    \item Implemented data structures and algorithms: Automata, Graph, Linked List, Stack, Queue, Heap Sort, Merge Sort, Quick Sort, Open address hash table with cluster diagnostics, Red Black Tree.
    \item C++14 conforming \texttt{stevemac::vector} with static allocator.
    \item Multithreaded, SIMD matrix multiplication.
    \item Rust code generator simple "wizard".
    \item FreeBSD/Linux Tree and List macros.
\end{itemize}

\subsection*{Algorithms Rust}
\begin{itemize}
    \item Implemented algorithms in Rust: Kadane, Linked List, Binary Tree, Quick Sort, Valid Parentheses, Binary Search, Insertion Sort, String Permute, Tic-tac-toe, etc.
\end{itemize}

\section*{Embedded Systems Project: ADC, DSP, and FFT Demonstration}
\subsection*{Project Overview}
The project involves an ADC (Analog-to-Digital Conversion), DSP (Digital Signal Processing), and FFT (Fast Fourier Transform) demo using custom data structures such as a priority queue and graph. This demonstration showcases advanced data structure usage in an embedded environment.

\subsection*{Steps and Basis for Edge Connections}
\begin{enumerate}
    \item \textbf{Priority Queue Construction}: 
    The FFT output is stored in a priority queue using \texttt{priority\_queue\_from\_array} and \texttt{priority\_queue\_build\_max\_heap}. The heap is then printed with \texttt{priority\_queue\_print\_heap}, organizing the FFT magnitudes based on their values (peaks and valleys).
    
    \item \textbf{Graph Creation}:
    A graph is created with \texttt{create\_graph(\&graph, FFT\_LENGTH)}, initializing the graph structure with a size based on the FFT length. Each vertex represents a magnitude from the FFT results.
    
    \item \textbf{Vertices Addition}:
    For each value in the heap (\texttt{pq.heap}), a vertex is added to the graph. Each vertex corresponds to an FFT output magnitude.
    
    \item \textbf{Edge Creation Based on Peaks and Valleys}:
    The core logic for adding edges checks for peaks and valleys in the magnitude values of the FFT results. Peaks are connected to neighboring valleys and vice versa.
    
    \item \textbf{Graph Visualization}:
    The graph is printed using \texttt{print\_graph(\&graph)}, showing the vertices (peaks and valleys) and their connections (edges).
\end{enumerate}

\subsection*{Edge Connection Summary}
\begin{itemize}
    \item Peaks are connected to adjacent valleys and vice versa.
    \item The edges are added based on relative magnitude comparisons of adjacent FFT results.
\end{itemize}

\section*{FreeBSD Kernel Deep Dive}
\begin{itemize}
    \item 45 hours of FreeBSD kernel study with McKusick’s deep dive course.
\end{itemize}

\section*{Rust Programming}
\begin{itemize}
    \item Completed the Rust book cover-to-cover, covering ownership, borrowing, smart pointers, concurrency, and more.
\end{itemize}

\section*{Embedded Systems and Development Tools}
\subsection*{Languages}
\begin{itemize}
    \item C, C++, ARM Assembly.
    \item Rust (embedded development).
\end{itemize}

\subsection*{Embedded Platforms}
\begin{itemize}
    \item ARM Cortex M series: STM32, Tiva C microcontrollers.
    \item Raspberry Pi.
\end{itemize}

\subsection*{Development Tools}
\begin{itemize}
    \item IAR Embedded Workbench, Keil uVision.
    \item STM32CubeMX (code generator).
    \item STM32CubeIDE, uVision Runtime Manager.
    \item Code Composer Studio, PinMux, Resource Explorer, Driverlib.
    \item Visual Studio, VisualGDB, VSCode, Vim, NeoVim.
    \item Rust embedded crates, ST-Link, OpenOCD.
    \item CLI environment (gdb + stlink or openocd).
    \item STM32 HAL, CMSIS Programming, Legacy StdPeripheral.
    \item MISRA Warnings enabled.
\end{itemize}

\subsection*{Embedded Programming Skills}
\begin{itemize}
    \item Advanced algorithms: Fast Fourier Transform, Automata, Finite State Machines.
    \item Electronic sampling: analog-to-digital, digital-to-analog, temperature, and voltage sampling.
    \item Data integrity techniques: wear leveling, CRC.
    \item Custom allocators using fixed array-based heap managers (e.g., Valvano).
    \item Standard C++ Library containers (e.g., \texttt{std::vector}) with fixed array-based custom allocators.
    \item Data structures optimized for embedded environments (stacks, queues, circular buffers, binary trees, priority queues, heaps, graphs, open-address hash tables).
    \item Algorithms derived from \textit{Introduction to Algorithms} by Cormen, Leiserson, Rivest, and Stein, Sedgewick’s \textit{Algorithms}, and \textit{Foundations of Computer Science} by Aho and Ullman.
    \item Automated embedded testing using CLI-based build, flash, and test sequences.
    \item Integration with STM32CubeMX for STM32 and PinMux/Resource Explorer for Tiva C series.
    \item CMSIS DSP Libraries for signal processing and testing.
\end{itemize}

\subsection*{Embedded Systems Skills}
\begin{itemize}
    \item General Operating Systems knowledge, Computer Organization, Binary Structure, and Embedded Startup Code Sequences.
    \item Timers and Clocks, generated programmatically or using tools like STM32CubeMX.
    \item Interrupt Handling, Fault Monitoring, and Recovery.
    \item Peripheral Communication: GPIO, UART, SSI, I2C, Timer, PWM, ADC, Ethernet, CAN.
    \item DMA streaming, Memory Management, Flash Memory, EEPROM, Power Optimization.
    \item Atomic Operations, Concurrency, and Parallelism (C++ Coroutines, Rust’s async/await).
    \item Debugging Tools: Keil uVision, STM32CubeIDE, Code Composer Studio, VisualGDB, CLI (gdb + st-link or openocd).
    \item Semihosting, UART Printing, ITM Tracing.
\end{itemize}

\end{document}
